\documentclass[a4paper,10pt]{article}
\usepackage[utf8]{inputenc}
\usepackage{graphicx}
\usepackage{subfigure}
\usepackage{fancyhdr}
\usepackage{listings}
\usepackage{ gensymb }


%opening

\pagestyle{fancy} %eigener Seitenstil
\fancyhf{} %alle Kopf- und Fußzeilenfelder bereinigen
\fancyhead[L]{SEP Abgabegespräch SS2015/16} %Kopfzeile links
\fancyhead[C]{} %zentrierte Kopfzeile
\fancyhead[R]{\includegraphics[width=1cm]{../../template/tugraz_bw}} %Kopfzeile rechts
\renewcommand{\headrulewidth}{0.4pt} %obere Trennlinie
\fancyfoot[C]{} %Seitennummer
\renewcommand{\footrulewidth}{0.4pt} %untere Trennlinie


\lstset{language=C++, basicstyle=\ttfamily\tiny}  

\begin{document}

\section*{Zug}

Schreibe ein einfaches Programm um die Passagiere von Zügen zu verwalten. Es können dabei mehrere Personen zu einem Zug hinzugefügt werden: Ein Schaffner, ein Lokführer und 'normale' Passagiere. 

\begin{description}
 \item[create $<$trainNumber$>$] Erstellt einen neuen Zug mit der angegebenen Nummer.
 \item[board $<$trainNumber$>$ $<$person$>$] Fügt eine Person zum Zug  $trainNumber$ hinzu. $person$ kann dabei die Werte 'pas' (passenger), 'con' (conductor)  und 'driver' (train driver) haben.
 \item[start $<$trainNumber$>$] Startet den Zug. Ein Zug kann nur gestartet werden, wenn mind. 1 Schaffner und 1. Lokführer im Zug sind. Falls der Zug gestartet wird, wird die Anzahl der Passagiere ausgegeben. 
\end{description}

Sollte bei \texttt{board} oder \texttt{start} die Zugnummer nicht existieren, soll ``Train does not exist!'' ausgegeben werden. Wenn bei $start$ noch ein Schaffner oder Lokführer fehlt, soll ``Cannot start train!``ausgegeben werden. Wenn bei \texttt{board} der Parameter nicht stimmt, soll "Cannot add <paramter>." ausgegeben werden.
Bei der Eingabe von \texttt{quit} oder EOF soll das Programm beendet werden.

\subsection*{Beispiel}
\begin{verbatim}
> create 234A
> board 234A pas
> board 234A con
> board 234A wrong
Cannot add wrong
> start 234A
Cannot start train!
> board 234A driver
> start 234A
Start train!
> board 234A driver
Train does not exist!
\end{verbatim}

\subsection*{Aufruf}
\texttt{./trains}

\subsection*{Bewertung}
Dokumentation \& Programmierstil,  Dokumentation \& Programmierstil, und Struktur \& Korrektheit (je 1 Punkt)
\newpage


\end{document}
